\chapter{Resultados}
\label{chap:resu}

En este capítulo se presentan los resultados experimentales del sistema de detección de intrusiones desplegado en hardware real, organizados en torno a las preguntas de investigación formuladas en el Capítulo~\ref{chap:Introduction}. Primero se analiza el desempeño de los modelos clásicos como línea base centralizada. Luego se comparan los modelos livianos entrenados en el notebook general. Posteriormente se evalúan las \textbf{ocho variantes HFL} \texttt{hfl\_v7-RN}, \texttt{hfl\_v7-no-ascon-RN}, \texttt{hfl\_v7-CNN} y \texttt{hfl\_v7-no-ascon-CNN}, cada una con sus modos \texttt{Results/} (NoFOG) y \texttt{Results\_FOG/}, integradas en una capa SRE unificada (\texttt{Analisis de Modelos/Completo/}) que cubre 58 intentos. Finalmente, se cuantifica el overhead del cifrado ASCON-128 contrastando las cuatro variantes con cifrado contra las cuatro PLAIN.

% ====================================================================
\section{Desempeño de Modelos Clásicos (Línea Base Centralizada)}
\label{sec:res_classic}

Los cuatro modelos clásicos de ML fueron entrenados sobre la totalidad del dataset (256\,276 flujos) con partición estratificada 80/20 y optimización de hiperparámetros mediante validación cruzada ($k=5$). La Tabla~\ref{tab:classic_results} resume los resultados obtenidos.

\begin{table}[H]
\centering
\caption{Desempeño de modelos clásicos de ML (entrenamiento centralizado).}
\label{tab:classic_results}
\small
\begin{tabular}{l c c c c c c}
\toprule
\textbf{Modelo} & \textbf{Accuracy} & \textbf{F1 Macro} & \textbf{F1 Weighted} & \textbf{Precision} & \textbf{Recall} & \textbf{MCC} \\
\midrule
Decision Tree      & 0.9987 & 0.9981 & 0.9987 & 0.9982 & 0.9981 & 0.9977 \\
Random Forest      & 0.9994 & 0.9991 & 0.9994 & 0.9991 & 0.9990 & 0.9989 \\
Logistic Regr.     & 0.9412 & 0.9118 & 0.9385 & 0.9244 & 0.9053 & 0.8964 \\
SVM (Linear)       & 0.9447 & 0.9162 & 0.9421 & 0.9280 & 0.9099 & 0.9017 \\
SVM (RBF)          & 0.9968 & 0.9953 & 0.9968 & 0.9954 & 0.9952 & 0.9944 \\
\bottomrule
\end{tabular}
\end{table}

\textbf{Hallazgos clave:}
\begin{itemize}
    \item \textbf{Random Forest} alcanza el mejor desempeño global (Accuracy: 99.94\%, MCC: 0.9989), seguido de cerca por Decision Tree (99.87\%) y SVM-RBF (99.68\%).
    \item Los modelos lineales (Logistic Regression, SVM-Linear) muestran una caída significativa ($\approx 5$--$6$ puntos porcentuales), lo que sugiere que las fronteras de decisión entre clases no son completamente linealmente separables en el espacio de 13 features.
    \item Todos los modelos basados en árboles presentan un MCC $> 0.99$, confirmando que las 13 features seleccionadas son altamente discriminativas para la tarea de clasificación triclase.
\end{itemize}

\subsection{Viabilidad de Despliegue en ESP32}

La Tabla~\ref{tab:classic_esp32} evalúa la factibilidad de exportar cada modelo a código C para ejecución en el ESP32-S3.

\begin{table}[H]
\centering
\caption{Estimación de huella de memoria para despliegue en ESP32.}
\label{tab:classic_esp32}
\small
\begin{tabular}{l r r l}
\toprule
\textbf{Modelo} & \textbf{Memoria estimada} & \textbf{Inferencia (ms)} & \textbf{Viable en ESP32?} \\
\midrule
Decision Tree       & $\sim$15\,KB  & $< 0.1$ & Sí \\
Random Forest       & $\sim$2.5\,MB & $\sim$5  & No (excede SRAM) \\
Logistic Regression & $\sim$1\,KB   & $< 0.05$ & Sí \\
SVM (Linear)        & $\sim$1\,KB   & $< 0.05$ & Sí \\
SVM (RBF)           & $\sim$4\,MB   & $\sim$50 & No (excede SRAM) \\
MLP (propuesto)     & $\sim$4.5\,KB & $< 0.1$  & \textbf{Sí} \\
\bottomrule
\end{tabular}
\end{table}

Esto revela un \emph{trade-off} fundamental: los modelos más precisos (Random Forest, SVM-RBF) son incompatibles con las restricciones de memoria del ESP32-S3 (512\,KB SRAM). El MLP propuesto ($\approx 4.5$\,KB) ofrece un compromiso eficiente entre precisión y viabilidad edge.

\section{Comparación de Modelos Livianos}
\label{sec:res_modelos_livianos}

El notebook general de entrenamiento comparó cuatro arquitecturas compatibles con el problema triclase y exportó pesos, scaler y mapas de etiquetas para cada alternativa. La Tabla~\ref{tab:light_models_results} resume los resultados offline observados.

\begin{table}[H]
\centering
\caption{Comparación offline de modelos livianos sobre las 13 features.}
\label{tab:light_models_results}
\small
\begin{tabular}{l c c l}
\toprule
\textbf{Modelo} & \textbf{Accuracy} & \textbf{F1 weighted} & \textbf{Uso en el proyecto} \\
\midrule
MLP/RN          & 0.9060 & 0.9049 & Rama principal \texttt{hfl\_v7-RN} \\
Residual-MLP    & 0.9055 & 0.9044 & Comparación offline \\
CNN-1D          & 0.9051 & 0.9041 & Rama experimental \texttt{hfl\_v7-CNN} \\
Tiny Transformer & 0.8903 & 0.8919 & Comparación offline; no desplegado \\
\bottomrule
\end{tabular}
\end{table}

Los resultados muestran que las tres arquitecturas compactas basadas en MLP/CNN se ubican alrededor de 90.5\% de accuracy offline. El MLP/RN se seleccionó como implementación principal porque obtiene el mejor resultado marginal y porque su inferencia en C++ es más simple que la CNN-1D. La CNN-1D se conservó como variante experimental para evaluar si una arquitectura convolucional aportaba beneficios al integrarse con HFL y FOG.

% ====================================================================
\section{Aprendizaje Federado Jerárquico: Convergencia y Precisión}
\label{sec:res_hfl}

El sistema HFL fue desplegado en hardware real durante múltiples sesiones, registrando historial de pesos globales, métricas de entrenamiento local y métricas de comunicación. Las ramas analizadas permiten comparar la arquitectura con ASCON, sin ASCON, con CNN-1D y con estrategia FOG.

\subsection{Convergencia del Modelo Federado}

\begin{table}[H]
\centering
\caption{Métricas de convergencia del sistema HFL.}
\label{tab:hfl_convergence}
\begin{tabular}{l r}
\toprule
\textbf{Métrica} & \textbf{Valor} \\
\midrule
Rondas hasta accuracy $> 70\%$  & $\sim$8--12 \\
Rondas hasta accuracy $> 85\%$  & $\sim$20--30 \\
Accuracy global estabilizada     & $\sim$0.90--0.95 \\
Loss final (Cross-Entropy)       & $\sim$0.14--0.25 \\
Muestras por ronda (por gateway) & 40 \\
Épocas locales por ronda         & 5 \\
Duración por ronda               & $\sim$51--61\,s según variante \\
\bottomrule
\end{tabular}
\end{table}

\textbf{Análisis:}
\begin{enumerate}
    \item El modelo federado converge en $\sim$20--30 rondas hacia un rango operativo de $\sim$90--95\% de accuracy, lo que representa un \textbf{gap} frente a los modelos centralizados más precisos. Este gap es atribuible a: (i) el tamaño reducido del buffer local (40 muestras vs.\ 256\,276), (ii) la distribución non-IID de los datos entre nodos, y (iii) la función heurística de etiquetado que introduce ruido respecto a las etiquetas reales.

    \item La convergencia mejora durante las primeras rondas y presenta oscilaciones posteriores, típicas de FL con datos non-IID \cite{Imteaj2022_FL_ResourceConstrainedSurvey}.

    \item El \textbf{costo de comunicación} por ronda federada completa (ida + vuelta) es de $\approx 7.2$\,KB en JSON plano y $\approx 9.8$\,KB con cifrado ASCON, lo cual es compatible con redes IoT de bajo ancho de banda.
\end{enumerate}

\subsection{Evolución temporal de accuracy, loss y pesos}
\label{subsec:curvas_convergencia}

Las Figuras~\ref{fig:hfl_v7_accuracy_loss_matrix} y~\ref{fig:hfl_v7_weight_magnitude_matrix} consolidan los resultados de las \textbf{ocho variantes} cubiertas por el análisis SRE unificado de \texttt{Analisis de Modelos/Completo/}: \texttt{RN/CNN} $\times$ \texttt{ASCON/PLAIN} $\times$ \texttt{NoFOG/FOG}. Para mantener homogeneidad visual entre corridas (algunas RN llegaron hasta ronda 30, otras quedaron en 20) las figuras se cortan en la \textbf{ronda 20}; las tablas SRE (\texttt{global\_round\_sli\_summary.csv}) conservan todas las rondas registradas. La primera figura compara la evolución conjunta de \textit{accuracy} y \textit{loss} global; la segunda resume la magnitud promedio de los pesos federados (\texttt{w3\_mag}, \texttt{w4\_normal}, \texttt{w4\_brute}, \texttt{w4\_scan}), útil para verificar si las actualizaciones se mantienen acotadas o si aparece deriva durante el entrenamiento.

\begin{figure}[H]
    \centering
    \includegraphics[width=0.98\textwidth]{images/Resultados/Completo/global_accuracy_loss.png}
    \caption{Accuracy y loss globales promedio por ronda (1--20) para las 8 variantes HFL v7. Generado por \texttt{sre\_completo\_analysis.py} a partir de los CSV de \texttt{global\_weights\_history\_*}.}
    \label{fig:hfl_v7_accuracy_loss_matrix}
\end{figure}

\begin{figure}[H]
    \centering
    \includegraphics[width=0.98\textwidth]{images/Resultados/Completo/weight_magnitude_trends.png}
    \caption{Magnitud promedio de los cuatro pesos federados (\texttt{w3\_mag}, \texttt{w4\_normal}, \texttt{w4\_brute}, \texttt{w4\_scan}) por variante a lo largo de las rondas 1--20.}
    \label{fig:hfl_v7_weight_magnitude_matrix}
\end{figure}

Estas curvas complementan la Tabla~\ref{tab:hfl_convergence}: el sistema no solo reporta métricas finales, sino trayectorias de entrenamiento comparables entre modelos y configuraciones. Tres lecturas clave de las Figuras~\ref{fig:hfl_v7_accuracy_loss_matrix}--\ref{fig:hfl_v7_weight_magnitude_matrix}:
\begin{enumerate}
    \item Las ocho variantes \textbf{convergen}: la accuracy global supera 0.90 entre las rondas 5 y 10 y la loss baja de $\sim$0.8 a $\sim$0.1--0.2 en el mismo rango.
    \item Las variantes con \textbf{CNN-1D} muestran una dinámica de pesos más expresiva: \texttt{CNN\_ASCON\_NoFOG} y \texttt{CNN\_ASCON\_FOG} aumentan progresivamente \texttt{w3\_mag} y los pesos asociados a las clases (\texttt{w4\_brute}, \texttt{w4\_scan}), mientras que las variantes RN se estabilizan rápidamente alrededor de \texttt{w3\_mag}~$\approx$~0.29.
    \item No se observa divergencia en ninguna variante dentro de las primeras 20 rondas. La rama \texttt{CNN\_PLAIN\_NoFOG} muestra el mayor ruido en \texttt{w4\_normal}, consistente con un entrenamiento más sensible a la inicialización local cuando no hay agregación FOG ni cifrado autenticado que filtre mensajes alterados.
\end{enumerate}

\subsection{Comparación: Federado vs.\ Centralizado}

La Tabla~\ref{tab:fed_vs_central} compara directamente la precisión del modelo federado con los modelos centralizados, considerando que el modelo federado utiliza un MLP idéntico.

\begin{table}[H]
\centering
\caption{Comparación de precisión: entrenamiento federado vs.\ centralizado.}
\label{tab:fed_vs_central}
\small
\begin{tabular}{l c c l}
\toprule
\textbf{Enfoque} & \textbf{Accuracy} & \textbf{Datos compartidos} & \textbf{Privacidad} \\
\midrule
MLP centralizado             & $\sim$97\%    & Todos los flujos      & Ninguna \\
MLP federado (HFL, 60 rondas) & $\sim$90\%   & Solo pesos (163 params) & Alta \\
Random Forest centralizado   & 99.94\%       & Todos los flujos      & Ninguna \\
Decision Tree centralizado   & 99.87\%       & Todos los flujos      & Ninguna \\
\bottomrule
\end{tabular}
\end{table}

El sistema federado alcanza $\sim$90\% de accuracy compartiendo \textbf{únicamente 163 parámetros del modelo} (652 bytes) en lugar de los datos crudos de tráfico ($\sim$256\,K flujos $\times$ 13 features $\approx$ 13.3\,M valores). Esto representa una \textbf{reducción superior al 99.99\%} en la cantidad de información intercambiada, con una penalización de $\sim$7--10 puntos en accuracy.

\subsection{Comparación entre variantes \texttt{hfl\_v7}}

La Tabla~\ref{tab:variants_results} consolida las \textbf{ocho combinaciones experimentales} cubiertas por el análisis SRE unificado de \texttt{Analisis de Modelos/Completo/}. Los valores son promedios sobre los intentos disponibles por variante y se calculan exclusivamente a partir de los CSV crudos generados durante las corridas (\texttt{ascon\_metrics\_*}, \texttt{plain\_metrics\_*}, \texttt{model\_metrics\_gateway\_*}, \texttt{global\_weights\_history\_*}). En total se procesaron \textbf{58 intentos}.

\begin{table}[H]
\centering
\caption{Resumen comparativo de las 8 variantes HFL v7 evaluadas. \textbf{GW p95}: latencia p95 de descifrado en gateway (\texttt{ESP32$\to$RPi}). \textbf{Srv dec p95}: latencia p95 de descifrado en servidor (\texttt{RPi$\to$PC}). \textbf{Overhead}: bytes adicionales promedio sobre el plaintext (0 en variantes \texttt{PLAIN}).}
\label{tab:variants_results}
\scriptsize
\begin{tabular}{l c c c c c c c}
\toprule
\textbf{Variante} & \textbf{Intentos} & \textbf{Acc.\ final} & \textbf{Loss final} & \textbf{Ronda (s)} & \textbf{GW p95 (ms)} & \textbf{Srv dec p95 (ms)} & \textbf{Overhead (B)} \\
\midrule
\texttt{CNN\_ASCON\_FOG}    &  7 & \textbf{0.9750} & \textbf{0.0910} & 58.61 & 4.53 & 37.47 & 1098 \\
\texttt{CNN\_ASCON\_NoFOG}  &  8 & 0.9604 & 0.1040 & 58.97 & 4.55 & 24.46 & 1049 \\
\texttt{CNN\_PLAIN\_FOG}    &  3 & 0.9444 & 0.1190 & 54.86 & 0.16 &  0.00 &    0 \\
\texttt{CNN\_PLAIN\_NoFOG}  &  3 & 0.9389 & 0.1363 & 52.02 & 0.15 &  0.00 &    0 \\
\texttt{RN\_ASCON\_FOG}     &  5 & 0.9633 & 0.1124 & 71.13 & 4.37 & 50.47 & 1084 \\
\texttt{RN\_ASCON\_NoFOG}   & 15 & 0.9381 & 0.1545 & 72.91 & 4.65 & 26.69 & 1050 \\
\texttt{RN\_PLAIN\_FOG}     &  3 & 0.9000 & 0.2199 & 67.14 & 0.16 &  0.00 &    0 \\
\texttt{RN\_PLAIN\_NoFOG}   & 14 & 0.9298 & 0.1581 & 51.01 & 0.17 &  0.00 &    0 \\
\midrule
\textbf{Promedio global}    & \textbf{58} & \textbf{0.9437} & \textbf{0.1383} & \textbf{61.36} & \textbf{2.34} & \textbf{17.39} & \textbf{636} \\
\bottomrule
\end{tabular}
\end{table}

La comparación de las 8 variantes confirma cinco hallazgos:
\begin{enumerate}
    \item \textbf{El ganador absoluto es \texttt{CNN\_ASCON\_FOG}} (acc 0.975, loss 0.091): combinar la mayor capacidad expresiva (CNN-1D) con la pre-agregación FOG y la integridad de ASCON produce el mejor punto de operación dentro del corpus.
    \item \textbf{El cifrado no degrada la convergencia}; al contrario, en este corpus las variantes \texttt{*\_ASCON\_*} tienden a igualar o superar a sus contrapartes \texttt{*\_PLAIN\_*}. Una explicación plausible es que la validación de tag autenticado \emph{filtra} mensajes corruptos o duplicados antes de que entren al buffer de entrenamiento.
    \item \textbf{FOG ayuda diferenciadamente} segun la familia de modelo y la seguridad: en CNN mejora $+1.5$ pts (ASCON) y $+0.6$ pts (PLAIN); en RN mejora $+2.5$ pts con ASCON pero \emph{degrada} $-3.0$ pts en PLAIN, lo que sugiere que la pre-agregación FOG es más sensible al ruido de mensajes sin autenticación cuando el modelo es lineal.
    \item \textbf{El costo de seguridad es asumible}: ASCON añade $\sim$1\,KB de overhead por mensaje y entre 4--5\,ms en el gateway (ESP32-S3 corre nativo en C++). En el servidor, las variantes FOG presentan latencias de descifrado un 50--100\% mayores que NoFOG por la cadena \texttt{peer$\to$leader$\to$PC}.
    \item \textbf{La duración de ronda en CNN es más estable} ($\sim$58\,s) que en RN ($\sim$51--73\,s); esta dispersión proviene de que las primeras corridas RN incluyen tiempos de calentamiento o reconexión Wi-Fi no presentes en las corridas CNN, no de un costo intrínseco mayor del modelo RN.
\end{enumerate}

La Figura~\ref{fig:hfl_v7_round_duration} muestra la duración por ronda promediada por variante, y la Figura~\ref{fig:hfl_v7_gateway_skew} ilustra el desbalance temporal entre gateway A y B medido como $|acc_A - acc_B|$ por ronda; el skew converge a valores bajos en todas las variantes después de la ronda 8--10.

\begin{figure}[H]
    \centering
    \includegraphics[width=0.95\textwidth]{images/Resultados/Completo/round_duration.png}
    \caption{Duración promedio por ronda federada para las 8 variantes (rondas 1--20).}
    \label{fig:hfl_v7_round_duration}
\end{figure}

\begin{figure}[H]
    \centering
    \includegraphics[width=0.95\textwidth]{images/Resultados/Completo/gateway_accuracy_skew.png}
    \caption{Desbalance de \emph{accuracy} entre gateways ($|acc_A - acc_B|$) por ronda y por variante.}
    \label{fig:hfl_v7_gateway_skew}
\end{figure}

\subsection{Comparación con la literatura}
\label{subsec:lit_comparison}

Para situar los resultados en el panorama actual, la Tabla~\ref{tab:literatura_comparativa} contrasta este sistema con trabajos representativos de FL+IDS y FL+TinyML. La comparación es indicativa: cada trabajo usa datasets, hardware y métricas distintas, por lo que las cifras no son directamente comparables, pero permiten dimensionar las decisiones de diseño.

\begin{table}[H]
\centering
\caption{Comparación cualitativa con trabajos relacionados.}
\label{tab:literatura_comparativa}
\scriptsize
\begin{tabular}{p{0.20\textwidth} p{0.10\textwidth} p{0.10\textwidth} p{0.08\textwidth} p{0.10\textwidth} p{0.10\textwidth} p{0.18\textwidth}}
\toprule
\textbf{Trabajo} & \textbf{Edge HW} & \textbf{Modelo} & \textbf{FL/HFL} & \textbf{Cifrado canal} & \textbf{Acc.\ reportada} & \textbf{Particularidad} \\
\midrule
Hajj~\textit{et~al.}~\cite{Hajj2023_CrossLayerFL_IDS} & PC sim. & MLP \emph{cross-layer} & FL plano & No & $\sim$0.97 & FL para IDS IoT con coordinación cross-layer \\
Chaurasia~\textit{et~al.}~\cite{Chaurasia2024_FL_ResNet_IIoT} & PC sim. & ResNet & FL plano & No & $\sim$0.99 & ResNet como cliente FL en IIoT \\
Ficco~\textit{et~al.}~\cite{Ficco2024_FL_TinyML_Onboard} & ESP32, RPi3, Arduino & NN compacta + TL & FL+TL on-board & No & 0.86 (clasif.) & FL+TL on-board, datasets ECG/Car \\
Roy~\textit{et~al.}~\cite{Roy2023_FL_AdaptedModel} & PC sim. & DNN & FL personalizado & No & $\sim$0.96 & Modelo localmente adaptado \\
Tran~\textit{et~al.}~\cite{Tran2026_EdgeTrustShard} & RPi Pico (264\,KB) & FL+blockchain & HFL+MF-PoP & Sí (cripto-PoP) & 0.94 a 30\% Byz. & Byzantine-robust FL en MCU \\
Heydari \& Mahmoud~\cite{HeydariMahmoud2025_TinyML_IDS} & MCU smart home & ML liviano & No FL & No & --- & TinyML IDS sin federación \\
Saravana Balaji~\textit{et~al.}~\cite{SaravanaBalaji2023_EdgeIoT_DL} & Edge sanitario & DBN+BOA & No FL & LABE & 0.979 & Edge salud con cifrado por atributos \\
\midrule
\textbf{Esta tesis (RN+ASCON)} & ESP32-S3 + RPi4 & MLP 1\,139 par. & HFL bidireccional & ASCON-128 & 0.946 & FL+TinyML+ASCON-128 medido \\
\textbf{Esta tesis (CNN+FOG)} & ESP32-S3 + RPi4 & CNN-1D & HFL+FOG & ASCON-128 & 0.975 & Pre-agregación FOG entre RPi \\
\bottomrule
\end{tabular}
\end{table}

Lecturas de la tabla:
\begin{itemize}
    \item Trabajos de FL para IDS sin restricciones edge \cite{Hajj2023_CrossLayerFL_IDS, Chaurasia2024_FL_ResNet_IIoT, Roy2023_FL_AdaptedModel} alcanzan 0.96--0.99, pero usan PCs como nodos cliente, sin demostrar despliegue en MCU. Esta tesis acepta una caída de $\sim$2--5 puntos a cambio de validar el ciclo en ESP32-S3 real con cifrado autenticado activo en todos los canales.
    \item Frente a Ficco~\textit{et~al.}~\cite{Ficco2024_FL_TinyML_Onboard}, que reporta 0.86 en clasificación, las variantes con ASCON y FOG superan ese punto sin sacrificar el ciclo on-board.
    \item Frente a EdgeTrust-Shard~\cite{Tran2026_EdgeTrustShard}, los resultados de accuracy son comparables aunque el problema atacado es distinto (Byzantine-robust FL vs.\ confidencialidad/integridad).
    \item Esta tesis es de los pocos trabajos identificados que combinan, en el mismo sistema desplegado y medido, HFL bidireccional + ASCON-128 + dos familias de modelos (MLP/CNN) + agregación FOG opcional.
\end{itemize}

\subsection{Limitaciones del estudio}
\label{subsec:limitations}

La interpretación de los resultados debe contemplar las siguientes limitaciones:
\begin{enumerate}
    \item \textbf{Tres clases y dataset acotado:} El problema se reduce a tres clases relevantes para tráfico MQTT (\texttt{normal}, \texttt{mqtt\_bruteforce}, \texttt{scan\_A}). Datasets más amplios como CICIoT2023, Edge-IIoTset y TON\_IoT \cite{Neto2023_CICIoT2023, Ferrag2022_EdgeIIoTset, Moustafa2020_TONIoT} contienen más vectores de ataque y permitirían generalizar mejor las conclusiones.
    \item \textbf{Etiquetado heurístico en gateway:} Las reglas que asignan etiquetas a las features recibidas son una aproximación calibrada con las distribuciones reales del dataset, pero introducen ruido. Esto contribuye al \emph{gap} de $\sim$7 puntos contra el modelo centralizado y es consistente con lo reportado en FL con datos non-IID \cite{Imteaj2022_FL_ResourceConstrainedSurvey, Albanbay2025}.
    \item \textbf{Número limitado de gateways y nodos:} Las pruebas usan $K=2$ gateways. Topologías con $K \geq 5$ y mayor heterogeneidad de cliente podrían exhibir convergencia más lenta o requerir agregadores más robustos como FedProx \cite{Li2020FedProx}.
    \item \textbf{Sin defensas contra clientes maliciosos:} Tal como se delimitó en la Sección~\ref{sec:threat_model}, el sistema no implementa Krum, FLTrust, Fed-NGA ni MF-PoP \cite{Tran2026_EdgeTrustShard}; un cliente comprometido puede sesgar la agregación.
    \item \textbf{Análisis de canal lateral pendiente:} La implementación de ASCON-128 no incluye contramedidas explícitas contra ataques de potencia, EM o \emph{fault injection} sobre el ESP32 \cite{Kaur2025_ASCON_Survey}.
    \item \textbf{Tráfico simulado en parte:} Aunque el ciclo se ejecuta sobre hardware real, el nodo \emph{simulated} reproduce distribuciones del dataset; un despliegue con captura de tráfico genuino aportaría mayor validez externa.
\end{enumerate}

% ====================================================================
\section{Overhead del Cifrado ASCON-128}
\label{sec:res_ascon}

Las métricas de cifrado/descifrado se recopilaron a lo largo de los \textbf{58 intentos} consolidados, abarcando todas las variantes \texttt{*\_ASCON\_*} y comparándose contra las correspondientes \texttt{*\_PLAIN\_*}. Esta sección consolida los hallazgos por canal y por dispositivo. Las figuras referenciadas se generan automáticamente en \texttt{Analisis de Modelos/Completo/} y se incluyen en \texttt{images/Resultados/Completo/}.

\begin{figure}[H]
    \centering
    \includegraphics[width=0.95\textwidth]{images/Resultados/Completo/transport_latency_p95.png}
    \caption{Latencia p95 (ms) promedio por variante y por canal/operación. Las variantes \texttt{*\_PLAIN\_*} solo muestran la latencia de serialización JSON (sub-milisegundo); las \texttt{*\_ASCON\_*} incluyen el costo de cifrado/descifrado autenticado.}
    \label{fig:hfl_v7_transport_p95}
\end{figure}

\begin{figure}[H]
    \centering
    \includegraphics[width=0.95\textwidth]{images/Resultados/Completo/transport_payload_bytes.png}
    \caption{Tamaño promedio de plaintext (izquierda) y overhead promedio sobre plaintext (derecha) por variante y canal. El overhead de las variantes \texttt{PLAIN} es 0 por construcción.}
    \label{fig:hfl_v7_transport_payload}
\end{figure}

\subsection{Overhead Temporal}

La Tabla~\ref{tab:ascon_time} resume los tiempos p95 promedio por variante para las tres operaciones que ocurren en cada ronda federada.

\begin{table}[H]
\centering
\caption{Latencia p95 (ms) promedio de cifrado/descifrado por variante, derivada de \texttt{transport\_sli\_summary.csv} en \texttt{Analisis de Modelos/Completo/}.}
\label{tab:ascon_time}
\small
\begin{tabular}{l c c c}
\toprule
\textbf{Variante} & \textbf{ESP32$\to$RPi (decrypt)} & \textbf{RPi$\to$PC (decrypt)} & \textbf{PC$\to$RPi (encrypt)} \\
\midrule
\texttt{CNN\_ASCON\_FOG}     & 4.53\,ms & 37.47\,ms & 40.37\,ms \\
\texttt{CNN\_ASCON\_NoFOG}   & 4.55\,ms & 24.46\,ms & 24.12\,ms \\
\texttt{RN\_ASCON\_FOG}      & 4.37\,ms & 50.47\,ms & 59.18\,ms \\
\texttt{RN\_ASCON\_NoFOG}    & 4.65\,ms & 26.69\,ms & 24.14\,ms \\
\midrule
\texttt{CNN\_PLAIN\_FOG}     & 0.16\,ms &  0.00\,ms &  0.48\,ms \\
\texttt{CNN\_PLAIN\_NoFOG}   & 0.15\,ms &  0.00\,ms &  0.62\,ms \\
\texttt{RN\_PLAIN\_FOG}      & 0.16\,ms &  0.00\,ms &  0.64\,ms \\
\texttt{RN\_PLAIN\_NoFOG}    & 0.17\,ms &  0.00\,ms &  0.25\,ms \\
\bottomrule
\end{tabular}
\end{table}

\textbf{Observaciones:}
\begin{itemize}
    \item \textbf{ESP32-S3 cifrando en C++}: la latencia p95 de descifrado en el gateway oscila entre $\approx$4.4--4.7\,ms en todas las variantes ASCON, prácticamente independiente del modelo (RN/CNN) y la topología (NoFOG/FOG). En las variantes \texttt{PLAIN} la latencia cae a $\approx$0.16\,ms (deserialización JSON pura).
    \item \textbf{Servidor PC en Python}: el descifrado en \texttt{RPi$\to$PC} ronda 24--27\,ms en NoFOG y sube a 37--50\,ms en FOG. La inflación FOG se debe a que el \emph{leader} reenvía mensajes preagregados con un tag adicional, lo que duplica las operaciones criptográficas en el camino.
    \item Para las variantes \texttt{*\_PLAIN\_*}, el descifrado del servidor reporta 0\,ms porque la deserialización JSON ocurre dentro del \texttt{request handler} de FastAPI sin instrumentación dedicada; la única fuente medida es el \texttt{encrypt} de despliegue (0.25--0.64\,ms).
    \item Asumiendo 3 operaciones criptográficas por ronda en el servidor, el overhead temporal de ASCON oscila entre $\sim$73\,ms (NoFOG) y $\sim$140\,ms (FOG). Sobre rondas de 51--73\,s, esto representa entre el \textbf{0.10\% y el 0.27\%} del tiempo total por ronda.
\end{itemize}

\subsection{Overhead de Comunicación}

\begin{table}[H]
\centering
\caption{Overhead promedio de tamaño por variante ASCON. \textbf{Plaintext}: bytes promedio de payload sin cifrar (\texttt{pt\_bytes}). \textbf{Overhead}: bytes adicionales del envelope ASCON+Base64 (\texttt{overhead\_bytes}). \textbf{Ratio}: porcentaje sobre el plaintext (\texttt{overhead\_bytes / pt\_bytes}). Las variantes \texttt{PLAIN} se omiten porque el overhead es 0 por construcción.}
\label{tab:ascon_size}
\small
\begin{tabular}{l r r r}
\toprule
\textbf{Variante} & \textbf{Plaintext (B)} & \textbf{Overhead (B)} & \textbf{Ratio (\%)} \\
\midrule
\texttt{CNN\_ASCON\_FOG}    & $\sim$2\,365 & 1\,098 & 46.46 \\
\texttt{CNN\_ASCON\_NoFOG}  & $\sim$2\,143 & 1\,049 & 48.94 \\
\texttt{RN\_ASCON\_FOG}     & $\sim$2\,323 & 1\,084 & 46.64 \\
\texttt{RN\_ASCON\_NoFOG}   & $\sim$2\,162 & 1\,050 & 48.55 \\
\midrule
\textbf{Promedio ASCON}     & $\sim$2\,248 & \textbf{1\,070} & \textbf{47.65} \\
\bottomrule
\end{tabular}
\end{table}

El overhead promedio del $\approx$48\% sobre el plaintext (en variantes ASCON) se compone de:
\begin{itemize}
    \item \textbf{Tag de autenticación ASCON:} 16 bytes (fijo).
    \item \textbf{Nonce:} 16 bytes (fijo).
    \item \textbf{Expansión Base64:} factor $\times 4/3$ sobre el ciphertext + tag + nonce.
    \item \textbf{Envelope JSON:} $\sim$30 bytes de estructura (\texttt{\{"ct":"...","tag":"...","nonce":"..."\}}).
\end{itemize}

En términos absolutos, el overhead es de $\approx$1.07\,KB por mensaje, lo que para las $\sim$3 transmisiones por ronda en NoFOG equivale a $\sim$3.2\,KB adicionales; en FOG, las dos saltos extra (peer$\to$leader, leader$\to$PC) elevan el overhead a $\sim$5--6\,KB por ronda. Para un enlace Wi-Fi estándar (54\,Mbps), esto añade $<$1\,ms de tiempo de transmisión, siendo prácticamente despreciable.

\subsection{Resumen: Impacto de ASCON en el Sistema HFL}

\begin{table}[H]
\centering
\caption{Resumen del overhead total de ASCON-128 por ronda federada, agregando los promedios de las 4 variantes ASCON y las 4 variantes PLAIN.}
\label{tab:ascon_resumen}
\begin{tabular}{l r r}
\toprule
\textbf{Componente} & \textbf{Con ASCON (4 variantes)} & \textbf{Sin ASCON (4 variantes)} \\
\midrule
Latencia p95 ESP32$\to$RPi (decrypt)   & 4.52\,ms      & 0.16\,ms \\
Latencia p95 RPi$\to$PC (decrypt)      & 34.77\,ms     & 0.00\,ms \\
Overhead promedio por mensaje          & $\sim$1.07\,KB & 0\,KB \\
Ratio de overhead sobre plaintext      & $\sim$48\%    & 0\% \\
Duración promedio de ronda             & 65.4\,s       & 56.3\,s \\
Accuracy global final promedio         & \textbf{0.9592} & 0.9272 \\
Loss global final promedio             & 0.1155        & 0.1583 \\
\midrule
\textbf{Overhead temporal relativo}    & \multicolumn{2}{c}{\textbf{0.10--0.27\% del ciclo FL}} \\
\textbf{Overhead de tamaño relativo}   & \multicolumn{2}{c}{\textbf{$\approx$48\% por mensaje}} \\
\textbf{Diferencia de accuracy}        & \multicolumn{2}{c}{\textbf{+3.2 pts a favor de ASCON}} \\
\bottomrule
\end{tabular}
\end{table}

Tres conclusiones de la tabla consolidada:
\begin{enumerate}
    \item ASCON-128 añade un overhead temporal \textbf{inferior al 0.3\% del ciclo federado} en cualquier configuración, y un overhead de tamaño moderado ($\sim$48\% del plaintext, $\sim$1\,KB absolutos por mensaje) compatible con redes IoT típicas.
    \item Las variantes \texttt{*\_ASCON\_*} \textbf{no solo no degradan la convergencia: la mejoran} en promedio $+3.2$ puntos de accuracy frente a las \texttt{*\_PLAIN\_*}. Esto refuerza la hipótesis de que la validación de tag autenticado actúa como filtro implícito de mensajes corruptos o duplicados.
    \item La diferencia se concentra en RN: \texttt{RN\_ASCON\_FOG} (0.9633) supera por 6.3 pts a \texttt{RN\_PLAIN\_FOG} (0.9000) y por 0.8 pts a \texttt{RN\_PLAIN\_NoFOG} (0.9298). En CNN la diferencia es menor (\texttt{CNN\_ASCON\_FOG} 0.9750 vs \texttt{CNN\_PLAIN\_FOG} 0.9444 = $+3.0$ pts).
\end{enumerate}

% ====================================================================
\section{Análisis de las Preguntas de Investigación}
\label{sec:preguntas}

\subsection{PI-1: ¿Puede TinyML detectar ataques de red en dispositivos IoT con precisión operativa?}

El modelo MLP de 1\,139 parámetros ($\sim$4.5\,KB), ejecutado directamente en el ESP32-S3, alcanza una accuracy operativa alrededor de 90--95\% tras el entrenamiento federado, dependiendo de la corrida y la variante. Si bien existe un gap frente al Random Forest centralizado (99.94\%), este último requiere $\sim$2.5\,MB de memoria y no es desplegable en un microcontrolador con recursos restringidos.

Comparando con modelos viables para edge (Decision Tree: 99.87\% con $\sim$15\,KB), el gap se reduce a $\sim$10 puntos, atribuible principalmente al esquema federado y no a la arquitectura del modelo. Un MLP idéntico entrenado de forma centralizada alcanza $\sim$97\%, situándose a solo 3 puntos del Decision Tree.

\textbf{Conclusión PI-1:} TinyML puede detectar ataques MQTT con precisión viable para alerta temprana en un factor de forma compatible con microcontroladores IoT. La penalización respecto a modelos centralizados se compensa con ventajas de privacidad, despliegue distribuido y seguridad de comunicación.

\subsection{PI-2: ¿Reduce el aprendizaje federado el intercambio de datos sensibles sin afectar la precisión del IDS?}

El aprendizaje federado reduce la información intercambiada de $\sim$13.3 millones de valores (todos los flujos) a \textbf{163 parámetros} por ronda (652\,bytes), una reducción de $>$99.99\%. El impacto en accuracy es de $\sim$7 puntos (de $\sim$97\% centralizado a $\sim$90\% federado), lo cual es consistente con la literatura de FL con datos non-IID \cite{Imteaj2022_FL_ResourceConstrainedSurvey, Albanbay2025}.

\textbf{Conclusión PI-2:} El aprendizaje federado elimina la necesidad de transferir datos crudos de tráfico, protegiendo la privacidad de las redes monitoreadas. La reducción de precisión ($\sim$7\%) es un trade-off aceptable en escenarios donde la privacidad es prioritaria.

\subsection{PI-3: ¿Cuál es el overhead del cifrado ASCON en la comunicación de modelos federados?}

Basado en el análisis SRE consolidado de 58 intentos sobre 8 variantes (Tabla~\ref{tab:variants_results}):
\begin{itemize}
    \item \textbf{Overhead temporal:} latencia p95 de descifrado ASCON de $\approx$4.5\,ms en el ESP32-S3 (C++) y $\approx$25--50\,ms en el servidor PC (Python). Sumando las 3--5 operaciones por ronda, el overhead total representa entre el \textbf{0.10\% y 0.27\%} del tiempo de ronda. En las variantes PLAIN, la latencia equivalente cae a sub-milisegundo (deserialización JSON pura).
    \item \textbf{Overhead de tamaño:} $\approx$48\% de incremento por mensaje (1.07\,KB sobre $\approx$2.25\,KB), atribuible al tag autenticado (16\,B), nonce (16\,B), expansión Base64 ($\times 4/3$) y envelope JSON.
    \item \textbf{Impacto en convergencia:} positivo, no negativo. Las variantes \texttt{*\_ASCON\_*} promedian \textbf{0.9592} de accuracy final frente a \textbf{0.9272} en las \texttt{*\_PLAIN\_*} ($+3.2$ pts). La hipótesis es que la validación obligatoria del tag autenticado descarta mensajes corruptos o duplicados antes de que entren al buffer de entrenamiento.
\end{itemize}

\textbf{Conclusión PI-3:} El overhead del cifrado ASCON-128 es \textbf{despreciable en tiempo} ($<0.3\%$ del ciclo FL) y \textbf{moderado en tamaño} ($\approx$48\%, $\sim$1\,KB absolutos), y además \textbf{mejora la accuracy} del modelo federado al filtrar implícitamente mensajes adulterados. Esto valida su idoneidad como mecanismo de protección end-to-end para HFL sobre IoT restringido.

\subsection{PI-4: ¿Qué aporta comparar MLP/RN, CNN-1D, baseline sin ASCON y FOG?}

La comparación de las 8 variantes (Tabla~\ref{tab:variants_results}) muestra que el proyecto no depende de un único script, sino de una arquitectura replicable donde se pueden intercambiar modelo, seguridad y topología, y que cada eje aporta una conclusión cuantificable:
\begin{itemize}
    \item \textbf{Modelo (RN vs CNN-1D):} CNN-1D supera consistentemente a RN cuando se mantiene el resto fijo. \texttt{CNN\_ASCON\_FOG} (0.9750) $>$ \texttt{RN\_ASCON\_FOG} (0.9633); \texttt{CNN\_ASCON\_NoFOG} (0.9604) $>$ \texttt{RN\_ASCON\_NoFOG} (0.9381). El costo es mayor complejidad de inferencia en C++ sobre ESP32, pero la accuracy ganada compensa en escenarios donde la disponibilidad de cómputo del MCU lo permita.
    \item \textbf{Seguridad (ASCON vs PLAIN):} ASCON añade $<$0.3\% de tiempo de ronda y $\sim$48\% de tamaño, pero \emph{mejora} la accuracy promedio en $+3.2$ pts. La rama PLAIN, lejos de ser solo un baseline, evidencia el valor de la validación de integridad como mecanismo \emph{de facto} de control de calidad de entradas.
    \item \textbf{Topología (NoFOG vs FOG):} FOG mejora la accuracy en CNN ($+1.5$ pts ASCON, $+0.6$ pts PLAIN) y en RN+ASCON ($+2.5$ pts), pero degrada \texttt{RN\_PLAIN\_FOG} ($-3.0$ pts vs \texttt{RN\_PLAIN\_NoFOG}). Esto sugiere que la pre-agregación FOG amplifica errores cuando el canal peer$\to$leader no tiene autenticación, y refuerza la conveniencia de exigir ASCON cuando se introduce la capa fog.
    \item \textbf{Modularidad operativa:} las cuatro carpetas \texttt{hfl\_v7-RN}, \texttt{hfl\_v7-no-ascon-RN}, \texttt{hfl\_v7-CNN}, \texttt{hfl\_v7-no-ascon-CNN}, cada una con \texttt{Results/} y \texttt{Results\_FOG/}, comparten esquema de CSV, lo que permitió construir una capa SRE única (\texttt{Analisis de Modelos/Completo/}) que produce métricas, eventos canónicos y trazas comparables entre todas las combinaciones sin reentrenar.
\end{itemize}

\textbf{Conclusión PI-4:} La arquitectura HFL v7 es modular y los cuatro ejes (modelo, seguridad, topología, escala) producen resultados cuantitativamente distintos. La combinación \textbf{CNN-1D + ASCON-128 + FOG} es la dominante en accuracy (0.975) y loss (0.091); cuando los recursos del MCU favorezcan al MLP, la mejor opción es \textbf{RN + ASCON + FOG} (0.963).
